
\documentclass{article}
\usepackage{colortbl}
\usepackage{makecell}
\usepackage{multirow}
\usepackage{supertabular}

\begin{document}

\newcounter{utterance}

\centering \large Interaction Transcript for game `cladder', experiment `full\_v1.5\_default', episode 771 with selfrefine{-}Qwen3.5{-}9B.
\vspace{24pt}

{ \footnotesize  \setcounter{utterance}{1}
\setlength{\tabcolsep}{0pt}
\begin{supertabular}{c@{$\;$}|p{.15\linewidth}@{}p{.15\linewidth}p{.15\linewidth}p{.15\linewidth}p{.15\linewidth}p{.15\linewidth}}
    \# & \multicolumn{2}{c}{Player} && \multicolumn{2}{c}{Game Master} \\
    \hline

    \theutterance \stepcounter{utterance}  
    & & & \multicolumn{4}{p{0.6\linewidth}}{
        \cellcolor[rgb]{0.9,0.9,0.9}{
            \makecell[{{p{\linewidth}}}]{
                \texttt{\tiny{[P1$\langle$GM]}}
                \texttt{You are an expert in causal inference. The following question is not a typical commonsense query, but rather a meticulously designed question created by a professor specializing in causal inference, intended to assess the students' mastery of the course content.} \\
\texttt{Imagine a self{-}contained, hypothetical world with only the following conditions, and without any unmentioned factors or causal relationships: Gender has a direct effect on occupation and salary. Skill has a direct effect on occupation and salary. Occupation has a direct effect on salary. Skill is unobserved. For individuals who are not male and blue{-}collar workers, the probability of high salary is 26\%. For individuals who are not male and white{-}collar workers, the probability of high salary is 61\%. For individuals who are male and blue{-}collar workers, the probability of high salary is 55\%. For individuals who are male and white{-}collar workers, the probability of high salary is 68\%. For individuals who are not male and with low skill levels, the probability of white{-}collar job is 40\%. For individuals who are not male and with high skill levels, the probability of white{-}collar job is 15\%. For individuals who are male and with low skill levels, the probability of white{-}collar job is 2\%. For individuals who are male and with high skill levels, the probability of white{-}collar job is 24\%. The overall probability of high skill level is 11\%. If we disregard the mediation effect through occupation, would gender negatively affect salary?} \\
\texttt{Start your answer with 'Yes' or 'No', followed by additional reasoning or evidence to support your explanation.} \\
            }
        }
    }
     \\ \\

    \theutterance \stepcounter{utterance}  
    & \multicolumn{4}{p{0.6\linewidth}}{
        \cellcolor[rgb]{0.561,0.737,0.561}{
            \makecell[{{p{\linewidth}}}]{
                \texttt{\tiny{[P1$\rangle$GM]}}
                \texttt{problem} \\
\texttt{No} \\
\\ 
\texttt{**Reasoning:**} \\
\\ 
\texttt{To determine if gender negatively affects salary when disregarding the mediation effect through occupation, we must isolate the **Direct Effect** of gender on salary. This corresponds to the causal path \$G \textbackslash{}to Y\$ while blocking the indirect path \$G \textbackslash{}to O \textbackslash{}to Y\$. In a Structural Causal Model (SCM), the direct effect is best estimated by comparing the outcome \$Y\$ for different values of \$G\$ while holding the mediator \$O\$ constant (i.e., calculating \$P(Y|G, O)\$).} \\
\\ 
\texttt{Let's examine the conditional probabilities provided for high salary (\$Y\$) given Gender (\$G\$) and Occupation (\$O\$):} \\
\texttt{1.  **Blue{-}collar workers:**} \\
\texttt{    *   \$P(Y=High | G=Male, O=Blue) = 55\textbackslash{}\%\$} \\
\texttt{    *   \$P(Y=High | G=Female, O=Blue) = 26\textbackslash{}\%\$} \\
\texttt{    *   Difference (Male {-} Female): \$+29\textbackslash{}\%\$} \\
\texttt{2.  **White{-}collar workers:**} \\
\texttt{    *   \$P(Y=High | G=Male, O=White) = 68\textbackslash{}\%\$} \\
\texttt{    *   \$P(Y=High | G=Female, O=White) = 61\textbackslash{}\%\$} \\
\texttt{    *   Difference (Male {-} Female): \$+7\textbackslash{}\%\$} \\
\\ 
\texttt{In both strata of occupation, holding occupation constant, males have a higher probability of a high salary than females. The direct effect of being male (relative to female) is positive in both groups (\$+29\textbackslash{}\%\$ and \$+7\textbackslash{}\%\$). Therefore, the direct effect of gender on salary is positive, not negative.} \\
\\ 
\texttt{The question asks if gender would *negatively* affect salary if we disregard the mediation. Since the direct effect is positive, the answer is No. The observed total effect might be influenced by the indirect path (where females are less likely to be in high{-}paying white{-}collar roles due to skill/occupation biases), but the specific question regarding the effect *disregarding* mediation points strictly to the direct causal link, which is positive.} \\
            }
        }
    }
    & & \\ \\

    \theutterance \stepcounter{utterance}  
    & & & \multicolumn{2}{p{0.3\linewidth}}{
        \cellcolor[rgb]{0.9,0.9,0.9}{
            \makecell[{{p{\linewidth}}}]{
                \texttt{\tiny{[GM$|$GM]}}
                \texttt{ParseError: The response doesn't start with one of the required word {[}'yes', 'no'{]}, but problem} \\
            }
        }
    }
    & & \\ \\

    \theutterance \stepcounter{utterance}  
    & & & \multicolumn{2}{p{0.3\linewidth}}{
        \cellcolor[rgb]{0.9,0.9,0.9}{
            \makecell[{{p{\linewidth}}}]{
                \texttt{\tiny{[GM$|$GM]}}
                \texttt{game\_result = ABORT} \\
            }
        }
    }
    & & \\ \\

\end{supertabular}
}

\end{document}
